\chapter{Differential-drive kinematics and odometry}
\label{ch:odometry}

Controlling one wheel does not yet describe how an entire robot moves.
A mobile robot combines the motion of several wheels, and its position
changes in a fixed reference frame. This chapter develops that connection
for a differential-drive robot, then uses the resulting model to estimate
position and orientation from wheel measurements.

\section{Geometry and assumptions}
A \emph{differential-drive robot} has two independently driven wheels on
a common axle line. Each wheel has its own motor; the axle line need not
be a single mechanical shaft. A passive caster can provide support without
setting the drive velocities. This arrangement allows the robot to travel
straight, follow a curve, or turn in place.

Choose the midpoint $P$ between the drive wheels as the origin of the robot
frame. Let $x_R$ point forward, $y_R$ point left, and $z_R$ point upward.
In the world frame, the midpoint has coordinates $(x,y)$ and the heading
$\theta$ is measured counterclockwise from $x_W$ to $x_R$. The planar pose is
\begin{equation}
 q=\begin{bmatrix}x\\y\\\theta\end{bmatrix}.
 \label{eq:l09-pose}
\end{equation}
It is also the state of the kinematic model used here. A model that includes
motor or body dynamics would generally require additional state variables.

\begin{figure}[htbp]
\centering\begin{tikzpicture}[scale=.95]
\draw[axis] (0,0)--(7,0) node[right]{$x_W$};
\draw[axis] (0,0)--(0,4.8) node[above]{$y_W$};
\draw[projection] (0,2)--(3,2)--(3,0);
\node[below] at (3,0) {$x$};\node[left] at (0,2) {$y$};
\begin{scope}[shift={(3,2)},rotate=30]
\draw[fill=gray!8,rounded corners] (-.8,-1.1) rectangle (1.5,1.1);
\draw[thick] (0,-1.3)--(0,1.3);
\foreach \s in {-1,1} {
 \draw[fill=gray!35] (-.55,1.3*\s-.13) rectangle (.55,1.3*\s+.13);
}
\draw[fill=white] (1.1,0) circle (.12);
\draw[axis,robotframe] (0,0)--(2.4,0) node[right]{$x_R$};
\draw[axis,robotframe] (0,0)--(0,2.2) node[above]{$y_R$};
\draw[axis,sensorframe] (.65,1.3)--(1.9,1.3) node[right]{$v_L$};
\draw[axis,sensorframe] (.65,-1.3)--(2.3,-1.3) node[right]{$v_R$};
\draw[<->] (-1.1,-1.3)--node[left]{$d$}(-1.1,1.3);
\end{scope}
\fill (3,2) circle (.05) node[below left]{$P$};
\draw[projection] (3,2)--(5,2);
\draw[axis] (4.5,2) arc (0:30:1.5);
\node at (4.75,2.28) {$\theta$};
\node[align=left] at (7,4.1) {Top view\\wheel radius $r$};
\end{tikzpicture}

\caption{A differential-drive robot viewed from above. The frame origin
is the axle midpoint. Wheel-center separation is $d$, and each wheel has
radius $r$. Both wheel velocities are positive in the robot's forward direction.}
\label{fig:l09-geometry}
\end{figure}

Assume that the robot moves on a plane, the chassis is rigid, the wheel
radii are equal to $r$, and the wheel-center separation is $d>0$. Both
wheels remain in contact with the ground and roll without longitudinal
or lateral slip. These are ideal kinematic constraints, not a claim that
physical friction is infinite. They are useful when the available traction
is sufficient for the motion being considered.

The wheel angular velocities $\omega_R$ and $\omega_L$ describe rotation
about the wheel axles. The robot angular velocity $\omega=\dot\theta$
describes yaw about the vertical axis. These are different rotations.
Choose the sign of each wheel coordinate so that positive wheel rotation
produces forward rolling. Then the signed wheel-center velocities along
$x_R$ are
\begin{equation}
 v_R=r\omega_R,\qquad v_L=r\omega_L.
 \label{eq:l09-wheel}
\end{equation}
Here the subscripts R and L mean right and left wheels; they do not denote
coordinate frames. Motor-side encoder rates must first be converted through
the transmission as in Chapter~\ref{ch:motors}.

\section{From wheel velocities to robot motion}
Let $v$ be the forward velocity of $P$. A point at lateral coordinate $y_R$
on the axle has forward velocity $v-\omega y_R$: the contribution from
rotation changes sign on opposite sides of the midpoint. Since the right
wheel is at $y_R=-d/2$ and the left wheel at $y_R=d/2$,
\[
 v_R=v+\frac d2\omega,\qquad
 v_L=v-\frac d2\omega.
\]
Adding and subtracting these two equations gives
\begin{equation}
 \boxed{v=\frac{v_R+v_L}{2}=\frac r2(\omega_R+\omega_L),
 \qquad \omega=\frac{v_R-v_L}{d}
 =\frac r d(\omega_R-\omega_L).}
 \label{eq:l09-forward}
\end{equation}
The midpoint velocity is the mean of the two wheel velocities. Their
difference determines the turning rate. For the same velocity difference,
a larger wheel separation produces a smaller turning rate.

Several limiting cases help check the signs. Equal wheel velocities give
$\omega=0$, so the robot travels straight. Equal and opposite wheel
velocities give $v=0$, so the midpoint remains stationary while the body
turns. With both wheels moving forward, a faster right wheel gives positive
$\omega$, so the robot turns left.

\begin{figure}[htbp]
\centering\begin{tikzpicture}[x=.95cm,y=.8cm]
\foreach \x/\title in {0/Straight,4/Turn in place,8/Curve left} {
 \begin{scope}[shift={(\x,0)}]
 \draw[thick] (-.7,0)--(.7,0);
 \draw[fill=gray!25] (-.85,-.4) rectangle (-.55,.4);
 \draw[fill=gray!25] (.55,-.4) rectangle (.85,.4);
 \fill (0,0) circle (.045);
 \node at (0,-1.3) {\title};
 \node[below] at (-.7,-.4) {L};\node[below] at (.7,-.4) {R};
 \end{scope}
}
\draw[axis,sensorframe] (-.7,.5)--(-.7,1.6);
\draw[axis,sensorframe] (.7,.5)--(.7,1.6);
\draw[axis,robotframe] (0,.2)--(0,2);
\draw[axis,sensorframe] (3,.1)--(3,-.9);
\draw[axis,sensorframe] (4.7,.5)--(4.7,1.5);
\draw[axis,robotframe] (4.4,.5) arc (0:260:.4);
\draw[axis,sensorframe] (7.3,.5)--(7.3,1);
\draw[axis,sensorframe] (8.7,.5)--(8.7,1.8);
\draw[axis,robotframe] (8,.2) .. controls (8,1.4) and (7.5,1.9) .. (7,2.1);
\end{tikzpicture}

\caption{Three motions viewed from above, with forward pointing upward.
Arrow lengths indicate wheel speeds. In-place turning uses opposite wheel
directions; curved travel uses unequal forward speeds.}
\end{figure}

\subsection{The inverse mapping}
A motion planner may specify desired $v$ and $\omega$, while the wheel
controllers require wheel angular velocities. Solving
Equation~\eqref{eq:l09-forward} in the other direction gives
\begin{equation}
 \boxed{\omega_R=\frac{2v+d\omega}{2r},\qquad
 \omega_L=\frac{2v-d\omega}{2r}.}
 \label{eq:l09-inverse}
\end{equation}
The parameters $r$ and $d$ describe the hardware; $v$ and $\omega$ describe
the requested motion. The resulting wheel commands must still lie within
the motors' speed and acceleration capabilities. For odometry, measured
wheel velocities are preferable to requested ones, since a command is not
a guarantee of what the robot actually does.

\subsection{Worked example: wheel and body velocities}
Suppose $r=0.05\,\mathrm m$, $d=0.30\,\mathrm m$,
$\omega_R=6\,\mathrm{rad/s}$, and $\omega_L=2\,\mathrm{rad/s}$.
Then
\[
 v_R=0.30\,\mathrm{m/s},\quad v_L=0.10\,\mathrm{m/s},\quad
 v=0.20\,\mathrm{m/s},\quad \omega=\frac23\,\mathrm{rad/s}.
\]
The robot moves forward and turns left. Substituting these body velocities
into Equation~\eqref{eq:l09-inverse} recovers $6$ and $2\,\mathrm{rad/s}$.

\section{The unicycle model}
Once wheel velocities have been mapped to $(v,\omega)$, the chassis motion
can be described by a single oriented point at $P$. This abstraction is
called the \emph{unicycle model}. It preserves forward motion and turning,
without retaining the two individual wheels in the pose equations.

The midpoint cannot move sideways instantaneously under the no-slip
assumption. Its velocity in robot coordinates is therefore $[v\;0]\trans$.
Rotating this vector into world coordinates gives
\begin{equation}
 \begin{bmatrix}\dot x\\\dot y\end{bmatrix}
 =\begin{bmatrix}\cos\theta&-\sin\theta\\
                  \sin\theta&\cos\theta\end{bmatrix}
  \begin{bmatrix}v\\0\end{bmatrix}.
\end{equation}
Together with $\dot\theta=\omega$, the model becomes
\begin{equation}
 \boxed{\dot x=v\cos\theta,\qquad
        \dot y=v\sin\theta,\qquad \dot\theta=\omega.}
 \label{eq:l09-unicycle}
\end{equation}
Equivalently, with input $u=[v\;\omega]\trans$,
\begin{equation}
 \dot q=
 \begin{bmatrix}\cos\theta&0\\\sin\theta&0\\0&1\end{bmatrix}u.
 \label{eq:l09-matrix}
\end{equation}
The state has three components, while the input has two. A dot denotes a
time derivative, so $\dot x$ and $\dot y$ have units of meters per second,
and $\dot\theta$ has units of radians per second.

\begin{figure}[htbp]
\centering\begin{tikzpicture}[scale=.9]
\draw[axis] (0,0)--(7,0) node[right]{$x_W$};
\draw[axis] (0,0)--(0,4.5) node[above]{$y_W$};
\coordinate (P) at (2,1.2);
\coordinate (V) at (5.5,3.5);
\draw[projection] (0,1.2)--(P)--(2,0);
\fill (P) circle (.05) node[below left]{$P$};
\draw[axis,sensorframe,very thick] (P)--node[above,sloped]{$v$}(V);
\draw[axis,robotframe] (P)--(.95,2.8) node[above]{$y_R$};
\node[right] at (V) {$x_R$};
\draw[projection] (V)--(5.5,1.2);
\draw[axis] (P)--node[below]{$v\cos\theta$}(5.5,1.2);
\draw[axis] (5.9,1.2)--node[right]{$v\sin\theta$}(5.9,3.5);
\draw (3,1.2) arc (0:33.3:1);
\node at (3.3,1.52) {$\theta$};
\end{tikzpicture}

\caption{The forward velocity expressed in world coordinates. Its components
are $v\cos\theta$ and $v\sin\theta$; a zero lateral velocity in the robot
frame does not imply a zero world-frame $y$ velocity.}
\end{figure}

\subsection{The lateral-motion constraint}
Projecting the world velocity onto the robot's left-pointing axis gives
\begin{equation}
 -\dot x\sin\theta+\dot y\cos\theta=0.
 \label{eq:l09-constraint}
\end{equation}
This is a \emph{nonholonomic} constraint: it restricts instantaneous
velocities and cannot be replaced by a constraint on pose alone. The robot
can still reach a location to its side by turning and driving; it cannot
translate directly sideways while holding its heading fixed.

The choice of representative point matters. A point located a distance $a$
ahead of the axle midpoint has robot-frame velocity $[v\;a\omega]\trans$.
Thus the simple zero-lateral-velocity statement applies to the chosen axle
midpoint, not to every point of a turning chassis.

\subsection{What the abstraction includes}
The unicycle equations are kinematic: they relate velocities to changes
in pose. They do not predict the forces, torques, balance, or motor transients
needed to achieve those velocities. They are useful when low-level control
and traction make the commanded motion a reasonable approximation, or when
measured velocities are used directly.

Similar abstractions can be useful for planning the paths of more complex
systems. Their admissible inputs must still reflect the actual platform.
For example, a car-like vehicle generally cannot turn in place, even though
an unconstrained unicycle model allows $v=0$ and $\omega\ne0$. A simplified
path model also does not replace a balance model for a walking robot.

\section{Odometry as integration of measured motion}
\emph{Wheel odometry} estimates changes in pose from wheel rotation.
Starting from an initial pose $q(t_0)$, the unicycle model implies
\begin{align}
 x(t)&=x(t_0)+\int_{t_0}^t v(\tau)\cos\theta(\tau)\,d\tau,\\
 y(t)&=y(t_0)+\int_{t_0}^t v(\tau)\sin\theta(\tau)\,d\tau,\\
 \theta(t)&=\theta(t_0)+\int_{t_0}^t\omega(\tau)\,d\tau.
\end{align}
This produces an estimate relative to the initial frame. Setting
$q(t_0)=[0\;0\;0]\trans$ places that frame at the starting pose; it does
not establish the robot's location in an independently defined building map.

On a computer, measurements arrive at discrete times $t_k$. Write
$h=t_{k+1}-t_k$ for the sampling interval and $q_k$ for the estimated pose
at $t_k$. The index $k+1$ means the next sample, not necessarily one second
later. If intervals vary, use their measured durations $h_k$.

\subsection{Forward Euler integration}
The derivative definition motivates the approximation
\[
 \dot x(t_k)\approx\frac{x(t_k+h)-x(t_k)}h.
\]
Using the velocity and heading at the beginning of the interval gives the
\emph{forward Euler} update:
\begin{equation}
 \boxed{\begin{aligned}
 x_{k+1}&=x_k+h v_k\cos\theta_k,\\
 y_{k+1}&=y_k+h v_k\sin\theta_k,\\
 \theta_{k+1}&=\theta_k+h\omega_k.
 \end{aligned}}
 \label{eq:l09-euler}
\end{equation}
Each new estimate becomes the starting point for the next step. All three
right-hand sides use the old state: updating the heading first and then
using that new heading for translation implements a different method.

For a general model $\dot q=f(q,u)$, the same construction is
\begin{equation}
 q_{k+1}=q_k+h f(q_k,u_k).
\end{equation}
This is numerical integration of a differential equation. The issue is
approximating motion over finite sampling intervals, not an inability of
computers to manipulate derivatives or limits symbolically.

\section{Which values belong in an integration step?}
Over one interval, the exact displacement depends on the velocity throughout
that interval. A numerical method specifies how the available information
approximates that motion. Choosing values consistently matters; arbitrary
mixtures of old and new variables need not have the intended accuracy.

For translation, the quantities being integrated are $v\cos\theta$ and
$v\sin\theta$. For heading, the integrand is $\omega$. The signed area
between an integrand and its approximation is the error in its integral.
A plot of heading alone does not directly give the error in position.

\begin{figure}[htbp]
\centering\begin{tikzpicture}[x=.85cm,y=.85cm]
\foreach \x/\title/\mode in {0/Left endpoint/0,4.5/Right endpoint/1,9/Trapezoid/2} {
 \begin{scope}[shift={(\x,0)}]
 \ifnum\mode=0
   \fill[robotframe!12] (.3,0) rectangle (2.8,.6);
   \draw[robotframe,thick,dashed] (.3,.6)--(2.8,.6);
 \fi
 \ifnum\mode=1
   \fill[sensorframe!12] (.3,0) rectangle (2.8,2.2);
   \draw[sensorframe,thick,dashed] (.3,2.2)--(2.8,2.2);
 \fi
 \ifnum\mode=2
   \fill[violet!12] (.3,0)--(.3,.6)--(2.8,2.2)--(2.8,0)--cycle;
   \draw[violet,thick,dashed] (.3,.6)--(2.8,2.2);
 \fi
 \draw[axis] (0,0)--(3.3,0) node[right]{$t$};
 \draw[axis] (0,0)--(0,2.7) node[above]{$g(t)$};
 \draw[thick] (.3,.6) .. controls (1.3,.7) and (1.8,2.1) .. (2.8,2.2);
 \draw[projection] (.3,0)--(.3,.6);\draw[projection] (2.8,0)--(2.8,2.2);
 \fill (.3,.6) circle (.045);\fill (2.8,2.2) circle (.045);
 \node[below] at (.3,0) {$t_k$};\node[below] at (2.8,0) {$t_{k+1}$};
 \node at (1.5,-.85) {\title};
 \end{scope}
}
\end{tikzpicture}

\caption{Approximating a scalar integrand $g(t)$ over one interval. Left
and right endpoint values give rectangles; a straight line between endpoint
values gives the trapezoidal rule. Its area equals the interval length times
the mean of those endpoint values.}
\label{fig:l09-quadrature}
\end{figure}

\subsection{Backward Euler}
Backward Euler evaluates the derivative at the new state and endpoint input:
\[
 q_{k+1}=q_k+h f(q_{k+1},u_{k+1}).
\]
In general this is an implicit equation that must be solved. For the
unicycle with externally supplied endpoint velocities, it can be evaluated
in sequence:
\begin{align*}
 \theta_{k+1}&=\theta_k+h\omega_{k+1},\\
 x_{k+1}&=x_k+h v_{k+1}\cos\theta_{k+1},\\
 y_{k+1}&=y_k+h v_{k+1}\sin\theta_{k+1}.
\end{align*}
Measured endpoint inputs are available only after that endpoint is reached.
This timing is compatible with reconstructing past motion, but it must not
be confused with predicting future motion from measurements not yet available.

\subsection{Trapezoidal integration}
The trapezoidal method averages the two endpoint \emph{derivatives}:
\[
 q_{k+1}=q_k+\frac h2
 \bigl[f(q_k,u_k)+f(q_{k+1},u_{k+1})\bigr].
\]
For the unicycle with known endpoint inputs, first compute
$\theta_{k+1}=\theta_k+\tfrac h2(\omega_k+\omega_{k+1})$, then use
\begin{align*}
 x_{k+1}&=x_k+\frac h2
  \bigl[v_k\cos\theta_k+v_{k+1}\cos\theta_{k+1}\bigr],\\
 y_{k+1}&=y_k+\frac h2
  \bigl[v_k\sin\theta_k+v_{k+1}\sin\theta_{k+1}\bigr].
\end{align*}
In particular,
\[
 \frac{v_k\cos\theta_k+v_{k+1}\cos\theta_{k+1}}2
 \ne
 \frac{v_k+v_{k+1}}2\cos\!\left(\frac{\theta_k+\theta_{k+1}}2\right)
\]
in general. Averaging the variables separately is a different approximation
because the velocity model is nonlinear in heading.

\subsection{Runge--Kutta methods}
Runge--Kutta denotes a family of methods with specified intermediate
derivative evaluations, rather than any operation that takes an average.
For example, the explicit second-order method known as Heun's method uses
an Euler prediction followed by an average of slopes:
\begin{align*}
 a_1&=f(q_k,u_k),&\widetilde q_{k+1}&=q_k+h a_1,\\
 a_2&=f(\widetilde q_{k+1},u_{k+1}),&
 q_{k+1}&=q_k+\frac h2(a_1+a_2).
\end{align*}
Here the second slope uses a \emph{predicted} endpoint state, whereas the
implicit trapezoidal method uses the final endpoint state. Inputs at the
required evaluation times must be supplied or approximated consistently.
Detailed solver design is beyond the present scope.\footnote{For the
standard Euler, Heun, and implicit trapezoidal formulas, see
\href{https://www.d.umn.edu/~mhampton/NA/ODEInt/}{Marshall Hampton's notes on numerical differential equations}.}

Both forward and backward Euler are first-order integration methods. The
trapezoidal and Heun methods are second order for sufficiently smooth
problems under the usual convergence assumptions. This refers to error
over a fixed time span as the step size decreases. A piecewise constant
approximation to an integrand should not be confused with calling Euler
a ``zeroth-order integration method.''

\subsection{Worked example: one interval, different methods}
Let $q_0=[0\;0\;0]\trans$, $h=1\,\mathrm s$, and suppose measured
endpoint velocities are
\[
 v_0=1\,\mathrm{m/s},\quad v_1=2\,\mathrm{m/s},\qquad
 \omega_0=0\,\mathrm{rad/s},\quad\omega_1=0.5\,\mathrm{rad/s}.
\]
Forward Euler uses the beginning values, so
$x_1=1$, $y_1=0$, and $\theta_1=0$. Backward Euler uses the endpoint
values and gives $\theta_1=0.5$, $x_1=2\cos(0.5)$, and
$y_1=2\sin(0.5)$. The trapezoidal heading is $\theta_1=0.25$, giving
$x_1=\tfrac12[1+2\cos(0.25)]$ and $y_1=\sin(0.25)$.

\begin{center}
\begin{tabular}{lrrr}
\toprule
Method & $x_1$ (m) & $y_1$ (m) & $\theta_1$ (rad)\\
\midrule
Forward Euler & 1.0000 & 0.0000 & 0.0000\\
Backward Euler & 1.7552 & 0.9589 & 0.5000\\
Trapezoidal & 1.4689 & 0.2474 & 0.2500\\
\bottomrule
\end{tabular}
\end{center}

The endpoint data alone do not uniquely determine the true trajectory
between samples. These are three approximations, not three exact answers.
If instead $u_0$ is explicitly held constant throughout the interval, the
exact motion is straight travel by one meter; the value $u_1$ then begins
the next interval. Distinguishing endpoint measurements from held commands
is part of specifying the problem.

\section{Sampling interval and missed motion}
Reducing the interval can improve the approximation when the motion changes
rapidly. Higher-order integration does not recover changes that were never
captured by the measurements. Figure~\ref{fig:l09-sampling} shows how an
apparently smooth interpolation can miss substantial variation.

\begin{figure}[htbp]
\centering\begin{tikzpicture}[x=1.6cm,y=1cm]
\draw[axis] (0,0)--(5.3,0) node[right]{$t$};
\draw[axis] (0,0)--(0,2.8) node[above]{$g(t)$};
\draw[thick] (0,.3) .. controls (.25,2.9) and (.6,2.9) .. (1,.5)
 .. controls (1.25,2.7) and (1.6,2.8) .. (2,.7)
 .. controls (2.3,3) and (2.65,2.5) .. (3,1)
 .. controls (3.4,.45) and (3.7,.45) .. (4,.4)
 .. controls (4.3,.3) and (4.7,.35) .. (5,.3);
\draw[violet,dashed,thick] (0,.3)--(1,.5)--(2,.7)--(3,1)--(4,.4)--(5,.3);
\foreach \x/\y in {0/.3,1/.5,2/.7,3/1,4/.4,5/.3} {
 \fill[violet] (\x,\y) circle (.045);
 \node[below] at (\x,0) {\x};
}
\end{tikzpicture}

\caption{Sparse samples can miss motion between sampling instants.
The solid curve represents a varying signal; the dashed line interpolates
the displayed samples. A more sophisticated solver cannot infer the hidden
peaks from these samples alone.}
\label{fig:l09-sampling}
\end{figure}

The useful interval depends on speed, turning rate, sensor timing, and
required accuracy. There is no universal sampling period for all ground
robots or all flying robots. Shorter windows also contain fewer encoder
counts and can make count-based speed estimates coarser, as discussed in
Chapter~\ref{ch:motors}. Processing cost, measurement resolution, and delay
must be considered together. Variable time steps are possible, provided
the measurement timestamps and integration intervals remain consistent.

\subsection{A constant-input benchmark}
For constant $v$ and $\omega\ne0$ over an interval, the unicycle equations
can be integrated exactly:
\begin{align}
 \theta_{k+1}&=\theta_k+h\omega,\\
 x_{k+1}&=x_k+\frac v\omega
       [\sin(\theta_k+h\omega)-\sin\theta_k],\\
 y_{k+1}&=y_k-\frac v\omega
       [\cos(\theta_k+h\omega)-\cos\theta_k].
 \label{eq:l09-exact}
\end{align}
For $\omega=0$, the motion is straight:
$x_{k+1}=x_k+hv\cos\theta_k$ and $y_{k+1}=y_k+hv\sin\theta_k$.
This special case provides a useful check on numerical integration, not
an exact model of arbitrary time-varying or slipping motion.

Starting at zero pose with $v=0.2\,\mathrm{m/s}$,
$\omega=0.5\,\mathrm{rad/s}$, and $h=1\,\mathrm s$, the exact endpoint is
approximately $(0.19177\,\mathrm m,0.04897\,\mathrm m,0.5\,\mathrm{rad})$.
A single forward Euler step gives $(0.2\,\mathrm m,0,0.5\,\mathrm{rad})$:
it updates heading but approximates the whole translation using the initial
direction. Smaller steps follow the curved path more closely.

\section{From wheel counts to an odometry update}
An encoder naturally measures rotation accumulated over an interval.
For signed counts $\Delta n_R,\Delta n_L$, motor-side count resolutions
$N_R,N_L$, and reductions $R_R,R_L$, the inferred wheel rotations are
\[
 \Delta\phi_R=\frac{2\pi\Delta n_R}{N_RR_R},\qquad
 \Delta\phi_L=\frac{2\pi\Delta n_L}{N_LR_L}.
\]
For the equal-radius model, define
\begin{equation}
 \Delta s=\frac r2(\Delta\phi_R+\Delta\phi_L),\qquad
 \Delta\theta=\frac r d(\Delta\phi_R-\Delta\phi_L).
\end{equation}
Then an Euler-style increment update is
\[
 x_{k+1}=x_k+\Delta s\cos\theta_k,\quad
 y_{k+1}=y_k+\Delta s\sin\theta_k,\quad
 \theta_{k+1}=\theta_k+\Delta\theta.
\]
The accumulated counts already describe the completed interval; they are
not instantaneous measurements of velocity at its end. With a short
interval, using $\theta_k+\Delta\theta/2$ for the translation direction
often improves the turning approximation. Under constant $v,\omega$, this
mid-heading approximation uses the arc length as a chord length; it is
still approximate. Total wheel increments alone do not determine the
translation exactly if the turning pattern varies within the interval.

A practical implementation initializes the pose and encoder baselines,
reads both count increments over matching intervals, converts counts to
wheel motion, and updates the pose using one consistent integration rule.
Retain the old pose until all components of the update have been computed.
It is convenient to integrate heading without wrapping and wrap it only
for display; directly averaging angles near $-\pi$ and $\pi$ can otherwise
give a misleading direction.

\clearpage
\section{Why odometry drifts}
Each update builds on the preceding estimate, so errors in motion and
heading carry into later updates. Sources include encoder quantization,
missed counts, incorrect wheel radii or separation, wheel slip, and numerical
integration error. A heading error is especially consequential because it
changes the direction used to project subsequent forward motion.

\begin{figure}[htbp]
\centering\begin{tikzpicture}[scale=.9]
\draw[axis] (0,0)--(8,0) node[right]{$x_W$};
\draw[axis] (0,0)--(0,4.4) node[above]{$y_W$};
\draw[thick] (.4,.4) .. controls (1,2) and (1.5,2.1) .. (2.4,1.6)
 .. controls (3.4,1) and (3.8,1.5) .. (4.2,2.5)
 .. controls (4.6,3.4) and (5.2,2.8) .. (6,3.3)
 .. controls (6.5,3.5) and (6.7,3.9) .. (7,4);
\draw[robotframe,dashed,thick] (.4,.4) .. controls (.9,2) and (1.3,2.7) .. (2,2.5)
 .. controls (2.8,2.3) and (3.3,2.8) .. (3,3.5)
 .. controls (2.8,3.8) and (3.1,4) .. (3.3,4.2);
\fill (.4,.4) circle (.05) node[below right]{start};
\node[right] at (7,4) {actual};
\node[robotframe,above] at (3.3,4.2) {odometry};
\end{tikzpicture}

\caption{Illustrative actual and estimated trajectories starting from the
same pose. Wheel odometry has no external position reference to pull the
estimate back toward the actual path. Separation need not increase at
every step, although uncertainty generally grows without correction.}
\end{figure}

Reducing the integration step addresses discretization error, but does not
remove slip or calibration bias. Conversely, drift is not inevitable solely
because integration is numerical: perfect motion information and an exact
model would recover the relative trajectory. In practice, wheel odometry
is useful over short intervals, while external observations such as satellite
positioning or recognized landmarks can constrain accumulated error.
Combining such information belongs to the broader problem of localization.

\clearpage
\section{Additional exercises}
Assume equal wheel radii, no slip, counterclockwise positive heading, and
positive wheel velocities for forward rolling. Use radians in computations.
\begin{enumerate}
\item \textbf{Wheel-to-body mapping.} Let $r=0.05\,\mathrm m$ and
$d=0.30\,\mathrm m$. Find $(v,\omega)$ for wheel rates
$(\omega_R,\omega_L)=(4,4),(4,-4),(6,2)$ in rad/s. Describe each motion
and check the units of both results.

\item \textbf{Inverse kinematics.} For the same robot, compute the wheel
rates required for $v=0.15\,\mathrm{m/s}$ and $\omega=0.4\,\mathrm{rad/s}$.
If each wheel is limited to $4\,\mathrm{rad/s}$ in magnitude, is the request
feasible? Scale both body commands by a common factor to fit the limit.

\item \textbf{World-frame velocity.} At heading $\theta=\pi/3$, the robot
has $v=0.2\,\mathrm{m/s}$ and $\omega=-0.5\,\mathrm{rad/s}$.
Compute $\dot q$ and verify Equation~\eqref{eq:l09-constraint}. Explain why
the robot can have nonzero $\dot y$ without moving sideways in its own frame.

\item \textbf{Two Euler steps.} Start at $q_0=[0\;0\;0]\trans$ with
$h=0.5\,\mathrm s$. Apply $(v_0,\omega_0)=(0.2,1)$ over the first interval
and $(v_1,\omega_1)=(0.2,0)$ over the second, in SI units. Compute $q_1,q_2$
with forward Euler. Explain why updating heading before position changes
the result.

\item \textbf{Integration accuracy.} For constant $v=0.2\,\mathrm{m/s}$
and $\omega=0.5\,\mathrm{rad/s}$, start at zero pose and integrate for one
second using one Euler step and then ten steps. Compare both position
estimates with Equation~\eqref{eq:l09-exact}. Which source of error is
being reduced, and which real-world errors are absent from this comparison?

\item \textbf{Encoder odometry.} Let $N_R=N_L=8$, $R_R=R_L=100$,
$r=0.05\,\mathrm m$, and $d=0.30\,\mathrm m$. Over one interval, record
40 right counts and 24 left counts. Find $\Delta\phi_R,\Delta\phi_L$,
$\Delta s$, and $\Delta\theta$. Update an initial pose
$(1\,\mathrm m,2\,\mathrm m,\pi/2)$ using the initial heading, then the
mid-heading approximation.

\item \textbf{Model and measurement limits.} Explain why decreasing the
step size cannot recover a brief wheel slip from encoder data alone. Give
an example of an external observation that could help correct pose drift.
Why is the final pose not determined uniquely by just two endpoint samples
of $v$ and $\omega$?
\end{enumerate}
